\chapter{Testing and results}
To verify the functionality of the developed firmware, the individual control loops were tested first after PI and PID controller tuning.
Afterwards, a practical usage example was demonstrated using CAN communication between a Python client and two uFOC drivers connected in a daisy-chain configuration.

For testing, two iPower GM4108 motors powered from a 20~V power supply were used.

\section{Control loops testing}
Testing of the individual control loops was performed by modifying the firmware main loop to periodically 
transmit the reference and measured variables over UART, while CAN communication was used to update the target values in real time.

\subsection{Torque loop testing}
The current control loop was tested by changing the target torque current from 0~A to 0.4~A, then to 0.75~A, 
and finally back to 0~A, while the motor shaft was physically blocked to allow the motor to generate maximum torque.

During the test, the target flux current was set to zero.

\begin{figure}[h]
    \centering
    \includegraphics[width=0.8\textwidth]{figs/current_testing.pdf}
    \caption{Torque control loop under load}
    \label{fig:current_test}
\end{figure}

Figure~\ref{fig:current_test} shows both torque and flux current components together with their target and measured values.
The measured torque current closely follows the commanded reference with minimal steady-state error and low noise.
At the same time, the flux current regulator successfully maintains the flux current near zero, 
with only brief transient deviations occurring during rapid changes of the torque current reference.

The results demonstrate stable operation of the implemented FOC current control loop 
and proper decoupling between the torque ($i_q$) and flux ($i_d$) current components.

\subsection{Velocity loop testing}

The velocity control loop was tested on an unloaded motor by changing the reference velocity from 0~rot/s to 4~rot/s, 
then to -8~rot/s to demonstrate operation in the opposite direction, and finally to 18~rot/s.

\begin{figure}[h]
    \centering
    \includegraphics[width=0.8\textwidth]{figs/velocity_testing.pdf}
    \caption{Velocity control}
    \label{fig:velocity_test}
\end{figure}

According to the manufacturer, the used GM4108 motor should achieve approximately 513--567~RPM, corresponding to 8.55--9.45~rot/s \cite{gm4108}.
However, testing showed that significantly higher rotational speeds could be achieved using the uFOC driver and the implemented FOC algorithm.

Figure~\ref{fig:velocity_test} shows that the controller is able to accurately track the reference velocity within the nominal operating range of the motor, 
including stable operation in both rotational directions.
When operating above the rated motor speed, the controller was still able to accelerate the motor further, reaching almost double the rated speed.

The results demonstrate correct functionality of the velocity controller and successful integration of the cascaded velocity and current control loops.


\subsection{Position loop testing}

The position control loop was tested on an unloaded motor by changing the target position from 0 to 3 rotations, then to 6 rotations, and finally back to 3 rotations.

\begin{figure}[h]
    \centering
    \includegraphics[width=0.8\textwidth]{figs/position_testing.pdf}
    \caption{Position control}
    \label{fig:position_test}
\end{figure}

Figure~\ref{fig:position_test} demonstrates smooth and stable position tracking of the implemented controller.
The measured position closely follows the target trajectory without significant steady-state error.
A non-zero derivative gain of the PID controller was used to damp the system response and reduce overshoot during rapid position changes.

The results confirm correct functionality of the cascaded position, velocity, and current control loops.

\section{Testing of client communication in daisy-chain mode}

To demonstrate a practical use case of the firmware, a simple mirror-mode example was implemented in \texttt{pytools/mirror\_mode.py}. 
The example uses two uFOC drivers, each connected to an iPower GM4108 motor.

The motors are connected in a daisy-chain CAN configuration. 
Each firmware instance is compiled with a different CAN ID, and one of the drivers is connected to a computer through a Makerbase CANable interface.

The Python script sets one motor to \texttt{NO\_CONTROL} mode and the other motor to \texttt{POSITION\_CONTROL} mode. 
In a continuous loop, the Python client reads the position of the first motor and sets the target position of the second motor to the same value. 
As a result, when the user manually moves the first motor, the second motor mirrors its motion.

This demonstration verifies the functionality of the CAN communication interface, runtime control mode switching, sensor feedback reading, 
and target position updating from the Python client. It also demonstrates that multiple uFOC drivers can share a single CAN bus in a daisy-chain configuration.

*add video reference*





